\subsection{Small indicators, large answers}
\label{sec:linearity-examples}

The difficult step is often choosing what to count. The following examples
use the same identity, $\E[\sum I_j]=\sum\P(I_j=1)$, on very different
objects. Each proof uses a finite probability space.

\begin{example}[Records in a random order]
Shuffle $n$ distinct scores uniformly and reveal them one at a time.
The $j$th score is a \emph{record} if it exceeds every previous score;
the first score is always a record. Among the first $j$ scores, symmetry
makes each position equally likely to hold the largest. Thus its record
indicator $I_j$ has mean $1/j$. The number $R_n$ of records satisfies
\[
\E[R_n]=\sum_{j=1}^n\frac1j=H_n,
\]
where $H_0=0$ and $H_n$ is the $n$th harmonic number. For example,
four distinct scores give mean $H_4=25/12$, although an individual
ordering has an integer number of records. We did not need to count
orderings with each possible record total.
\end{example}

\begin{example}[Inversions in a shuffle]
In a uniform permutation $\pi$ of $1,\ldots,n$, an \emph{inversion}
is a pair $i<j$ with $\pi(i)>\pi(j)$. For each of the $\binom n2$
pairs, swapping the two values pairs the greater-first and smaller-first
cases, so the inversion probability is $1/2$. Therefore
\[
\E[\text{number of inversions}]=\frac12\binom n2=\frac{n(n-1)}4.
\]
Indicators sharing a position need not be independent; linearity still
applies. For $n=1$ the empty sum is zero.
\end{example}

\subsection{Spaghetti loops: count the closures}
\label{sec:spaghetti-loops}

Start with $n\ge1$ separate strands, each with two distinguishable ends.
Pair all $2n$ ends uniformly at random and tie each pair together. There
are $(2n-1)(2n-3)\cdots1$ perfect matchings: choose the partner of a
specified end, remove those two ends, and repeat. Equivalently, repeatedly
choose any specified free end and choose its partner uniformly among
the other free ends. Every full matching has the same product probability.
Closed loops stay closed; only \emph{free} ends can be selected.

At any intermediate stage, every unfinished component is an open chain
with exactly two free ends. If there are $m$ open chains, there are
$2m$ free ends. Tying two ends has just two possibilities:

\begin{center}
\begin{tikzpicture}[thick,scale=.9]
\begin{scope}
\draw[mint,very thick] (-1,0) to[bend left=40] (1,0);
\draw[rose,dashed,very thick] (-1,0) to[bend right=40] (1,0);
\fill (-1,0) circle (2pt); \fill (1,0) circle (2pt);
\node[align=center,text width=5cm] at (0,-1)
{Same chain:\\one loop closes};
\end{scope}
\begin{scope}[xshift=7cm]
\draw[mint,very thick] (-1.5,0)--(-.5,0);
\draw[mint,very thick] (.5,0)--(1.5,0);
\draw[rose,dashed,very thick] (-.5,0)--(.5,0);
\foreach \x in {-1.5,-.5,.5,1.5}{\fill (\x,0) circle (2pt);}
\node[align=center,text width=5cm] at (0,-1)
{Different chains:\\one longer chain remains};
\end{scope}
\end{tikzpicture}
\end{center}

Either operation reduces the number of open chains by one. Given any
completed history with $m$ chains remaining, the selected end has
$2m-1$ possible partners, exactly one of which is the other end of
its own chain. The chance that this tie closes a loop is therefore
$1/(2m-1)$, regardless of the history.

Let $J_m$ indicate closure at the tie made with $m$ chains remaining.
Since the conditional probability has the same value for every possible
history, total probability gives $\E[J_m]=1/(2m-1)$.
Every final loop was closed by exactly one tie, so the loop count $L_n$ is
$\sum_{m=1}^nJ_m$. Hence
\[
\boxed{\E[L_n]=\sum_{m=1}^n\frac1{2m-1}
=H_{2n}-\frac12H_n.}
\]
For $n=1,2,3$ the means are $1,4/3,23/15$, respectively. With no
strands there are no loops, consistent with the empty sum at $n=0$.
The indicators refer to \emph{closure steps}, not to original strands:
one loop may contain several original strands. No independence assertion
is needed for the expectation calculation. This classical example also
appears in the Stanford EE 178/278A expectation notes listed in the references.

\subsection{An average proves that a good cut exists}
\label{sec:random-cut}

A finite graph consists of vertices and edges joining distinct vertices.
Suppose it has $e$ edges. Color each vertex red or blue independently
with equal probabilities. An edge crosses the resulting \emph{cut} if
its ends have different colors. For each edge this probability is
$2(1/2)(1/2)=1/2$, so
\[
\E[\text{number of crossing edges}]=e/2.
\]
At least one coloring has at least $e/2$ crossing edges: otherwise every
value, and hence their average, would be smaller than $e/2$.
Thus every such graph has a cut containing at least half its edges.
On a triangle the edge indicators are not mutually independent, since
the number of crossing edges can only be zero or two. Their dependence
does not obstruct the average argument.

\begin{bigidea}[title={Choose what contributes one}]
For records, count successful positions. For inversions, count pairs.
For loops, count closing ties. For a cut, count crossing edges.
The right decomposition can make a complicated distribution unnecessary.
\end{bigidea}
